\chapter*{Preface}
\addcontentsline{toc}{chapter}{Preface}
% STATUS: written 2026-11-06 (states the scope, the source rule and the living
% edition rule; these are the book's own policies, fixed in CLAUDE.md and in the
% ROADMAP, not external claims.)
% claim-ok markers added 2026-11-12: this preface states the policies of the book
% itself, which is what the marker exists for. The two external claims (Leibniz
% 1684, Ax 1971) carry their citations.
\label{page:preface}
% The scope statement below is what the origins and relevance judges read as the
% book's idea. A book has no abstract environment, so the idea is declared here,
% in an environment the judges can find, and the preface prose follows it.
\begin{abstract}
\noindent
A derivation obeys the product rule: an additive map $d$ with
$d(xy)=d(x)y+xd(y)$. Leibniz wrote that rule in 1684 for the calculus, and three
centuries later the same equation defines a map on any ring. This book follows the
equation from the calculus to the present, class by class. The classes begin
with
commutative rings, and continue with the prime and semiprime rings and the
Jordan, Lie, Leibniz, alternative, Malcev, Poisson, super and $n$-ary
algebras. They continue with path, incidence and
Leavitt path algebras; matrix, triangular, polynomial and Weyl
algebras; Banach and operator algebras; group and Hopf algebras; and
differential
fields with their Galois theory. The book follows the equation into the
maps that
weaken it as well, which are higher derivations, biderivations and other
multiadditive maps. They are the
local and 2-local derivations, the generalized and twisted derivations, and the
commuting and centralizing maps that a derivation must be.
The book
follows it too into the groups
that measure it. These are the Hochschild cohomology of an algebra and the
Gerstenhaber bracket on it, whose first cohomology group is the derivations
modulo the inner ones. For each class it
states the definitions with their located sources, records the theorems that
settle which derivations are inner, and reports the questions still open. Every
statement belongs to a source read in full from the page, and the book states no
result of its own.
\end{abstract}

A derivation is the algebraic form of a derivative \cite[p.~1]{Ritt1950}.
Leibniz wrote the product rule in 1684 \cite{Leibniz1684}, and three centuries later the same rule
defines a map on any ring: an additive map $d$ with $d(xy)=d(x)y+xd(y)$.
Every chapter of this book studies some version of that equation, either the % claim-ok: the plan of this book
equation itself or a weaker or stronger form of it.

\section*{Why an encyclopaedia and not a survey}

A survey selects and argues. An encyclopaedia tries to be complete in scope, so that a % claim-ok: what this book is
reader can find where a given question is treated and what is known about it.

This book is exhaustive in scope and brief in text. % claim-ok: what this book is
Every flavour of derivation, every class of ring and algebra where derivations % claim-ok: the scope of this book
act, and every line of development from Leibniz to this year has a place in the
plan.
The plan is fixed in advance and the text grows into it. % claim-ok: how this book is made

\section*{The source rule}

A book of derivations rests on the equation $d(xy)=d(x)y+xd(y)$ and on three % claim-ok: the premise of this book
centuries of work about it. Almost every statement in it therefore belongs to someone else. % claim-ok: a statement about this book
A reader is entitled to know whose it is, on which page it stands, and in which
of the nine languages the search was made. The rule below is how that is kept. % claim-ok: a statement about this book

No fact, definition, example or application is stated before its source has been % claim-ok: the rule this book follows
read in full from the page.
A recalled theorem appears as a theorem-like environment with its located % claim-ok: the rule this book follows
citation, keeping the kind and the wording of the source.
Proofs and remarks of the book itself are marked as its own. % claim-ok: the rule this book follows

This rule has a cost, and the cost is visible. % claim-ok: a statement about this book
A section whose sources are not yet obtained stays a marked planned % claim-ok: the rule this book follows
stub, and the stub says which source it waits for.
A reader who finds a stub knows that the coverage is planned there. % claim-ok: the rule this book follows
The text is not yet written, which is a different thing from a section that was % claim-ok: a statement about this book
written and then cut.

This edition has no such stub. % claim-ok: the state of this edition
Every section is written. % claim-ok: the state of this edition
Some sources cited in the bibliography could not be opened at all. % claim-ok: the state of this edition
Their entries carry the marker `unverified-ref` together with the reason. % claim-ok: the rule this book follows
Their sections cite them for what the record states, rather than quoting a
theorem from a page nobody has seen.
The count is kept in the download record of this book, not here, because it % claim-ok: where this book keeps its record
changes as copies are found.
The work list that the search step keeps records the same fact for the works % claim-ok: how this book is made
that were found and not cited. A lead is not a source until its page has been
read.
The policy stays in place for later editions, which is why it is stated here. % claim-ok: the rule this book follows

\section*{Where the applications are}

A derivation is a rule that obeys the Leibniz product rule. The rule itself
goes back to Leibniz \cite{Leibniz1684}. The name was already familiar to
Ritt when he built differential algebra on it \cite[p.~1]{Ritt1950}. The same
rule appears again in fields beyond the differential calculus, and each one is % claim-ok: the plan of this book
reported in the chapter that owns its notions, beside the definition it uses.

Quantum mechanics is in Chapter~\ref{ch:origins}, where Born, Jordan and Dirac % claim-ok: the plan of this book
write the commutator as a derivation. Coding theory is in
Chapter~\ref{ch:higher-multi}, where a Hasse--Schmidt derivation gives the
linear recurrence a linear code is built from. Differential geometry is in
Chapters~\ref{ch:polynomial-weyl} and~\ref{ch:commutative}, with vector fields
and the modules of differential operators over them. Group analysis, model
theory and differential Galois theory are in
Chapter~\ref{ch:differential-algebra}. Locally nilpotent derivations and the
Jacobian conjecture are in Chapter~\ref{ch:commutative}.

A reader who wants the applications of one notion finds them in the section % claim-ok: the arrangement of this book
that states it. Nothing is collected into a chapter of its own.
A reader who has to leave the object to find its application must come back to
the object afterwards. The two halves are then harder to hold together than the
one page they were on.

Physics asked early why a derivation arises. Dirac % claim-ok: a pointer to the section that states and cites this
\cite{Dirac1925} translated noncommutative kinematics into the language of a
Poisson bracket
(Section~\ref{sec:basic-theory:physics}), and Kadison \cite{Kadison1966} later
asked whether every derivation of an operator algebra comes from inside
(Section~\ref{sec:banach-operator:kadison-sakai}). Geometry came next: a tangent vector
at a point of a smooth manifold is a derivation of the smooth functions there
(Section~\ref{sec:commutative:geometry}). Differential algebra began when Ritt
and Kolchin axiomatized the symbol $d/dx$ as a derivation on a field, and the
field-of-definition question for the Painlev\'e equations sits inside that
axiomatization (Section~\ref{sec:differential-algebra:differential-equations}). In
positive characteristic the ordinary product rule survives in the
divided-power form of Hasse and Schmidt; the higher derivations are the
algebraic backbone of Reed--Solomon and Goppa codes
(Section~\ref{sec:higher-multi:coding}). Control theory discovered that an
input--output system is a finitely generated differential field extension,
with the input a differential transcendence basis
(Section~\ref{sec:differential-algebra:control}). On a group the same rule is a
derivation of the group ring, and Fox showed the free group carries one
derivation per generator (Section~\ref{sec:group-hopf:harmonic-analysis}).
Ax translated the conjecture of Schanuel about the complex exponential into a % claim-ok: a pointer to the section that states and cites this
statement about a derivation on a differential field \cite{Ax1971}. His 1971 theorem is the
seed of the Ax--Schanuel inequalities
(Section~\ref{sec:differential-algebra:transcendence}). Hochschild and cyclic
homology compute derivations as the homology of a chain complex, and on a
smooth commutative algebra the answer is the exterior algebra of K\"ahler
differentials (Section~\ref{sec:cohomology:algebraic-topology}). A computer
program that differentiates a function runs the chain rule and the Leibniz
rule step by step under automatic differentiation
(Section~\ref{sec:differential-algebra:automatic-differentiation}).

\hbfigure{02-basic-theory/oscillator}{The harmonic oscillator. The derivation $f\mapsto\{f,H\}$ with $H=(p^2+q^2)/2$ is the vector field $(p,-q)$ on the phase plane; its flow is a rotation.}{fig:basic:oscillator}

\section*{The living edition}

The subject grows. % claim-ok: a statement about this book's field
Papers on derivations appear every month, and a printed encyclopaedia would be out of % claim-ok: why this book is a living edition
date before it reached a reader.
This edition is therefore updated monthly: new work is added where it changes a % claim-ok: how this edition is maintained
statement, a table entry or an open problem.
Appendix~\ref{app:atlas} reports what the edition has read, and the version
number on the title page names the edition.

\section*{How the book was made}

A book of derivations is a book of equations, and an equation in it is worth
what its source is worth. Every statement about a derivation here was therefore
traced to a page before it was written. The work of reading, checking and
drafting was carried out with artificial % claim-ok: how this book was made
intelligence under human direction.
A person set the scope, fixed the questions and decided what the book would % claim-ok: how this book was made
contain.
Machine tools assisted the reading of each source, the extraction of its results % claim-ok: how this book was made
and the drafting of each section.
Every claim was then checked against the page it cites. % claim-ok: how this book was made
The rules that govern this are stated in the repository of the book. % claim-ok: where this book's rules live
The records of what was read are kept there too, so a reader can trace any % claim-ok: where this book's records live
statement to its source.

The name on the title page is the person who directed the work and answers for
it.
